\section{Síntesis y ampliación}

Esta clase parte de la pregunta «¿por qué los robots todavía no saben ordenar una casa?» y llega a una descomposición clara del sistema: el foundation model aporta el semantic prior, los robot data aportan el action grounding, la interface determina el nivel de abstracción de la salida del modelo, el controller se encarga de la dynamics y de las constraints, la evaluation verifica la transfer y la safety layer limita las consecuencias físicas. PaLM-E demuestra que el sensor embedding y el language model pueden entrenarse en conjunto; RT-2 conecta los action tokens a un VLM y forma un VLA; RT-X pone a prueba la cross-embodiment positive transfer; Language to Rewards, por su parte, deja que el LLM escriba el objective y que el MPC escriba la trajectory.

\subsection{Seis lecciones de ingeniería transferibles}

\begin{enumerate}
  \item \textbf{Primero se define la interfaz y después se elige el modelo.} Decidir si la salida es una skill, un action token, una trajectory, una reward o una constraint es más fundamental que aumentar los parámetros sin criterio.
  \item \textbf{Tratar la data mixture como objeto de diseño.} Los web data, los robot success, los failure, los recovery y los datos de varias embodiments requieren pesos y auditorías distintos.
  \item \textbf{Separar la semantic generalization de la motor generalization.} Reconocer un objeto nuevo no equivale a producir una acción de contacto nueva; las métricas de evaluación deben desglosarse.
  \item \textbf{Demostrar la ruta de la transfer con ablations.} Observar solo que el modelo grande funciona mejor no permite distinguir el efecto de la scale, el co-fine-tuning, los robot data y el web prior.
  \item \textbf{Conservar la visibilidad del controller y de la safety layer.} Aunque la policy sea de extremo a extremo, la capa de ejecución debe seguir teniendo guardrails de rate/force/geometry y un emergency stop.
  \item \textbf{Considerar el hybrid bootstrap.} El reward optimizer puede generar diverse data y el VLA puede someterse a distill para obtener una policy rápida; ambos pueden formar un data flywheel.
\end{enumerate}

\begin{importantbox}{Esta clase condensada en una frase}
La contribución más importante de los foundation models a la robótica no es que el modelo de lenguaje controle todo directamente, sino que nos permite rediseñar \textbf{las representaciones compartidas, los datos de entrenamiento y las interfaces de tarea}; un embodied system realmente utilizable todavía debe conectar semántica, acción, dinámica, evaluación y seguridad en un ciclo cerrado.
\end{importantbox}

\subsection{Las tres afirmaciones que más se malinterpretan}

\begin{warningbox}{No convertir evidencia parcial en conclusiones generales}
«RT-2 tiene emergent skills» no significa que haya aprendido cualquier acción nueva; «RT-X muestra positive transfer» no significa que la robotics scaling law ya esté madura; «el LLM puede escribir la reward» no significa que el código generado pueda llevarse directamente a un robot real. Las tres afirmaciones deben acompañarse de las condiciones de action support, dataset boundary, controller, validation y safety.
\end{warningbox}

\subsection{Lecturas adicionales}

\begin{itemize}
  \item Driess et al., \href{https://arxiv.org/abs/2303.03378}{\emph{PaLM-E: An Embodied Multimodal Language Model}}
  \item Brohan et al., \href{https://arxiv.org/abs/2212.06817}{\emph{RT-1: Robotics Transformer for Real-World Control at Scale}}
  \item Brohan et al., \href{https://arxiv.org/abs/2307.15818}{\emph{RT-2: Vision-Language-Action Models Transfer Web Knowledge to Robotic Control}}
  \item Open X-Embodiment Collaboration, \href{https://arxiv.org/abs/2310.08864}{\emph{Open X-Embodiment: Robotic Learning Datasets and RT-X Models}}
  \item Yu et al., \href{https://arxiv.org/abs/2306.08647}{\emph{Language to Rewards for Robotic Skill Synthesis}}
  \item Lynch et al., \href{https://arxiv.org/abs/2206.04717}{\emph{Interactive Language: Talking to Robots in Real Time}}
\end{itemize}

Al leer estos artículos puede seguirse usando la lista de verificación de cinco casillas de esta clase: observation, model, interface, controller, evidence. En cuanto alguna de ellas quede oculta tras las palabras «de extremo a extremo» o «emergent», hay que seguir preguntando de dónde vienen los datos, cómo se expresan las acciones, quién se encarga de las restricciones físicas, cómo se registran las fallas y qué experimento podría refutar la explicación actual.

\end{document}
